\chapter{Flight-Ready LiDAR-Inertial Odometry}
\label{ch:flight-ready}

\section{Introduction}
\label{sec:fr-introduction}

An odometry system that achieves high accuracy on benchmark datasets is not
necessarily suitable for closed-loop flight control. Flight controllers impose
specific requirements that benchmarks rarely evaluate:

\begin{enumerate}
    \item Pose feedback at hundreds of hertz for stable control loops.
    \item Linear and angular velocity estimates at the same rate, required for
    damping and feedforward control.
    \item Low-noise signals that do not excite oscillations in the controller.
    \item Uninterrupted continuity, even during transient sensor degradation.
\end{enumerate}

FAST-LIO2~\cite{xu2022fastlio2} publishes odometry only upon completing an
\gls{ieskf} update, which occurs at the accumulated \gls{lidar} arrival rate. On
the target platform, with the Livox Mid-360 operating at 100\,Hz and accumulating
10 scans per cycle, this rate is approximately 10\,Hz. At a flight speed of
2\,m/s, the controller receives a new pose every 0.2\,m and must extrapolate
between updates. Furthermore, the original odometry message does not populate the
\texttt{twist} field with velocities, forcing the flight controller to numerically
differentiate the pose---a process that amplifies noise and introduces an inherent
delay of at least one sampling period.

This chapter presents the architectural modifications that raise the odometry
publishing rate to the IMU rate (200\,Hz on the target platform), including
filtered linear and angular velocities in the body frame. The contributions are:

\begin{itemize}
    \item A forward propagation scheme that publishes pose and velocity at every
    IMU sample, using the last \gls{ieskf}-corrected state as an anchor.
    \item The derivation of linear and angular velocities in the body frame,
    absent from the original implementation.
    \item A first-order low-pass filter for velocity and orientation, with an
    analysis of the smoothness--lag trade-off.
    \item A multi-threaded execution architecture that guarantees a constant
    200\,Hz rate regardless of \gls{ieskf} processing load.
    \item A chained propagation mechanism that keeps odometry alive during
    transient \gls{lidar} interruptions.
\end{itemize}

\section{Limitations of the Original Architecture}
\label{sec:fr-limitations}

The publishing architecture of FAST-LIO2 exhibits three limitations that make it
unsuitable for flight control:

\subsection{Publishing Rate Coupled to LiDAR Rate}
\label{sec:fr-coupling}

In the original implementation, odometry publishing occurs inside the scan
processing loop. Each cycle consists of: accumulating a complete scan, executing
IMU propagation, performing the \gls{ieskf} iterations with point-to-plane
registration, updating the ikd-Tree, and finally publishing the pose. The
publishing rate is therefore coupled to the \gls{lidar} arrival rate, typically
10--20\,Hz.

This coupling has an additional consequence: when the processing of a scan exceeds
the inter-scan period (for example, during the insertion of many new points into
the ikd-Tree), publishing is delayed and the effective rate degrades below 10\,Hz.

\subsection{Absence of Velocities}
\label{sec:fr-no-velocity}

The ROS~2 \texttt{nav\_msgs::Odometry} message contains two main fields:
\texttt{pose} (position and orientation) and \texttt{twist} (linear and angular
velocities). FAST-LIO2 only populates the \texttt{pose} field, leaving the
\texttt{twist} field at zero. This forces the flight controller to numerically
differentiate the position to obtain velocity---a process that amplifies noise and
introduces an inherent delay of at least one sampling period.

However, the filter already estimates the velocity in the world frame as part of
the state vector (equation~\ref{eq:state-vector}). The gyroscope, corrected by the
estimated bias, directly provides the angular velocity in the body frame. Both
quantities are available without additional computation and are more accurate than
any numerical differentiation.

\subsection{Single-Threaded Execution}
\label{sec:fr-single-thread}

The original implementation uses a single ROS~2 executor
(\texttt{SingleThreadedExecutor}) that serializes all callbacks: IMU, \gls{lidar},
and \gls{ieskf} processing. While the filter processes a scan, the IMU callback is
blocked. This blocking is harmless when odometry is published only at the end of
processing, but becomes a fundamental obstacle when publishing at IMU rate, since
the IMU callback must execute without delay at every sample.

\section{Forward Propagation at IMU Rate}
\label{sec:fr-forward-prop}

\subsection{Operating Principle}
\label{sec:fr-principle}

The \gls{ieskf} correction occurs at the accumulated \gls{lidar} rate
($\sim$10\,Hz). Between corrections, each new IMU sample allows the state to be
propagated forward using the inertial dynamics equations. This propagation does not
replace the filter: it uses the last corrected state as a starting point (referred
to as the \emph{anchor}) and applies a single-step integration with the most recent
IMU measurement.

Let $t_a$ denote the time of the last \gls{ieskf} correction and
$(\mathbf{R}_a, \mathbf{p}_a, \mathbf{v}_a, \mathbf{b}_{g,a}, \mathbf{b}_{a,a},
\mathbf{g}_a)$ the corrected state at that instant. For an IMU sample at time
$t > t_a$, with measurements $\boldsymbol{\omega}_m$ and $\mathbf{a}_m$, the
bias-corrected quantities are defined as:

\begin{align}
    \boldsymbol{\omega} &= \boldsymbol{\omega}_m - \mathbf{b}_{g,a}
    \label{eq:fr-omega} \\
    \mathbf{a}_b &= \mathbf{a}_m - \mathbf{b}_{a,a}
    \label{eq:fr-abody}
\end{align}

The acceleration in the world frame is obtained by rotating the body-frame
acceleration and adding gravity:

\begin{equation}
    \mathbf{a}_w = \mathbf{R}_a \, \mathbf{a}_b + \mathbf{g}_a
    \label{eq:fr-aworld}
\end{equation}

With $\Delta t = t - t_a$, the propagated state is:

\begin{align}
    \mathbf{p}(t) &= \mathbf{p}_a + \mathbf{v}_a \, \Delta t
        + \tfrac{1}{2} \, \mathbf{a}_w \, \Delta t^2
    \label{eq:fr-prop-pos} \\
    \mathbf{v}(t) &= \mathbf{v}_a + \mathbf{a}_w \, \Delta t
    \label{eq:fr-prop-vel} \\
    \mathbf{R}(t) &= \mathbf{R}_a \, \mathrm{Exp}(\boldsymbol{\omega} \, \Delta t)
    \label{eq:fr-prop-rot}
\end{align}

where $\mathrm{Exp}: \mathbb{R}^3 \to \mathrm{SO}(3)$ is the exponential map that
converts a rotation vector into a rotation matrix via the Rodrigues formula.

\begin{remark}
This integration assumes that the anchor rotation $\mathbf{R}_a$, the biases
$\mathbf{b}_{g,a}, \mathbf{b}_{a,a}$, and gravity $\mathbf{g}_a$ remain constant
between \gls{ieskf} corrections. For $\Delta t < 0.1$\,s (the nominal case at
10\,Hz), the error of this approximation is on the order of centimeters in position
and fractions of a degree in orientation~\cite{sola2017quaternion}, negligible
relative to the flight controller's resolution.
\end{remark}

\subsection{Anchor Update}
\label{sec:fr-anchor}

Each time the \gls{ieskf} completes a correction with a new \gls{lidar} scan, the
corrected state is copied to the propagation anchor under mutual exclusion. The
anchor contains: position, velocity, rotation matrix, gyroscope and accelerometer
biases, gravity vector, and the timestamp of the scan end.

The next IMU sample takes the new anchor as its starting point, resetting
$\Delta t$ to a value close to zero. In this way, the propagation never accumulates
error beyond the interval between two consecutive \gls{ieskf} corrections.

\section{Body-Frame Velocities}
\label{sec:fr-velocities}

\subsection{Linear Velocity}
\label{sec:fr-linear-vel}

The \gls{ieskf} state vector includes velocity in the world frame,
$\mathbf{v} \in \mathbb{R}^3$. The flight controller requires velocity in the body
frame, which is obtained through the inverse rotation:

\begin{equation}
    \mathbf{v}_b(t) = \mathbf{R}(t)^\top \, \mathbf{v}(t)
    \label{eq:fr-vel-body}
\end{equation}

where $\mathbf{R}(t)$ and $\mathbf{v}(t)$ are the propagated values from
equations~\eqref{eq:fr-prop-rot} and~\eqref{eq:fr-prop-vel}, respectively.

Body-frame velocity has a practical advantage over world-frame velocity: its
components align with the vehicle axes (forward, lateral, vertical), allowing
direct use in the flight controller's velocity control loops.

\subsection{Angular Velocity}
\label{sec:fr-angular-vel}

The angular velocity in the body frame is obtained directly from the gyroscope
measurement corrected by the estimated bias:

\begin{equation}
    \boldsymbol{\omega}_b(t) = \boldsymbol{\omega}_m(t) - \mathbf{b}_{g,a}
    \label{eq:fr-omega-body}
\end{equation}

Unlike linear velocity, which requires propagation and integration, angular
velocity is a direct gyroscope measurement. The bias $\mathbf{b}_{g,a}$ is updated
at every \gls{ieskf} correction, ensuring that the correction is consistent with the
most recent state estimate.

\subsection{Absence in the Original Implementation}
\label{sec:fr-vel-absence}

In the original FAST-LIO2, the odometry publishing function constructs the message
with the pose (position and quaternion) but does not populate the velocity fields.
The \texttt{twist} field of the \texttt{nav\_msgs::Odometry} message remains at
zero. This is adequate for offline mapping applications, where velocity is not
needed, but insufficient for flight control.

The velocity $\mathbf{v}$ exists in the filter state vector and is updated during
both propagation and \gls{ieskf} correction. The angular velocity is available at
every IMU sample. Publishing these quantities requires no additional computation
beyond the frame transformation of equation~\eqref{eq:fr-vel-body}.

\section{Velocity and Orientation Filtering}
\label{sec:fr-filtering}

\subsection{Need for Filtering}
\label{sec:fr-filter-need}

The propagated velocity is computed by integrating the accelerometer from the last
corrected state. Accelerometer noise transfers directly to the propagated velocity,
producing high-frequency fluctuations on the order of 0.2\,m/s. A finely tuned
velocity controller would amplify these fluctuations, generating oscillations in the
actuators.

The propagated orientation suffers from an analogous problem: integrated gyroscope
noise produces fluctuations in the output quaternion, particularly visible in the
yaw angle.

\subsection{First-Order Low-Pass Filter}
\label{sec:fr-ema}

An exponential moving average (EMA) filter is applied to the published velocities.
Let $\alpha \in (0, 1]$ be the smoothing parameter. For the body-frame linear
velocity:

\begin{equation}
    \bar{\mathbf{v}}_b[k] = \alpha \, \mathbf{v}_b[k]
        + (1 - \alpha) \, \bar{\mathbf{v}}_b[k-1]
    \label{eq:fr-ema-vel}
\end{equation}

where $k$ indexes the IMU samples and $\bar{\mathbf{v}}_b[k]$ denotes the filtered
velocity. The same equation is applied to the angular velocity:

\begin{equation}
    \bar{\boldsymbol{\omega}}_b[k] = \alpha \, \boldsymbol{\omega}_b[k]
        + (1 - \alpha) \, \bar{\boldsymbol{\omega}}_b[k-1]
    \label{eq:fr-ema-omega}
\end{equation}

The cutoff frequency of the EMA filter at a sampling rate $f_s$ is:

\begin{equation}
    f_c = \frac{f_s}{2\pi} \cdot \ln\!\left(\frac{1}{1 - \alpha}\right)
    \label{eq:fr-cutoff}
\end{equation}

For $f_s = 200$\,Hz and $\alpha = 0.05$, the cutoff frequency is
$f_c \approx 1.6$\,Hz, which removes accelerometer noise without significantly
attenuating drone maneuvers, whose typical bandwidth is below 1\,Hz at moderate
speeds.

\subsection{Orientation Filtering via SLERP}
\label{sec:fr-slerp}

Quaternion filtering cannot be performed component-wise, as the weighted sum of two
unit quaternions does not in general produce a unit quaternion, and subsequent
normalization introduces distortions when the angular change is large.

Spherical linear interpolation (SLERP) respects the manifold structure of $S^3$
and provides the geodesic interpolation between two orientations. The orientation
filter is defined as:

\begin{equation}
    \bar{\mathbf{q}}[k] = \mathrm{SLERP}\!\left(\bar{\mathbf{q}}[k-1],\;
        \mathbf{q}[k],\; \alpha\right)
    \label{eq:fr-slerp}
\end{equation}

where $\mathbf{q}[k]$ is the propagated quaternion at sample $k$. For small
angular changes between consecutive samples (as is the case at 200\,Hz), SLERP is
practically equivalent to linear interpolation with normalization, but it is the
mathematically correct choice and avoids artifacts during aggressive maneuvers.

\subsection{Filtering in the Body Frame}
\label{sec:fr-body-filtering}

Linear velocity filtering is applied after the transformation to the body frame
(equation~\eqref{eq:fr-vel-body}), not before. If the world-frame velocity were
filtered and then rotated by $\mathbf{R}(t)$, the multiplication by the rotation
matrix would inject gyroscope noise into the filtered signal. By filtering
$\mathbf{v}_b$ directly, the rotation is applied once and the filter handles all
noise sources jointly.

\subsection{Smoothness--Lag Trade-off}
\label{sec:fr-tradeoff}

The parameter $\alpha$ controls the trade-off between noise reduction and signal
lag:

\begin{itemize}
    \item $\alpha = 1.0$: no filtering, direct propagator output.
    \item $\alpha = 0.1$: moderate smoothing, cutoff frequency $\approx 3.2$\,Hz.
    \item $\alpha = 0.05$: strong smoothing, cutoff frequency $\approx 1.6$\,Hz,
    lag of approximately 0.2\,s.
\end{itemize}

For flight applications with a maximum speed of 2\,m/s, $\alpha = 0.05$ provides a
reduction in velocity root-mean-square error of one order of magnitude compared to
the unfiltered signal, with acceptable lag for the controller bandwidth.

\section{Multi-Threaded Execution Architecture}
\label{sec:fr-multithreading}

\subsection{Problem: Rate Degradation under Load}
\label{sec:fr-degradation}

With a single-threaded executor (the default ROS~2 configuration), all callbacks
are serialized on the same thread. When the \gls{ieskf} processes a \gls{lidar}
scan, the IMU callback is queued until processing finishes. Since the processing
time of a scan can exceed 50\,ms (especially during massive insertions into the
ikd-Tree), the odometry publishing rate degrades from 200\,Hz to values as low as
44\,Hz.

This problem does not exist in the original implementation because odometry is
published inside the scan processing loop. But by moving publishing to the IMU
callback to achieve 200\,Hz, blocking by \gls{ieskf} processing becomes the main
bottleneck.

\subsection{Solution: IMU Executor Isolation}
\label{sec:fr-isolation}

The solution consists of using two independent single-threaded executors:

\begin{itemize}
    \item \textbf{IMU executor}: runs exclusively the IMU callback on a dedicated
    thread. Never blocked by scan processing.
    \item \textbf{Main executor}: runs the \gls{lidar} callbacks, the \gls{ieskf}
    timer callback, and the point cloud publishers.
\end{itemize}

The IMU callback is registered to an exclusive callback group
(\texttt{CallbackGroup}), which is assigned to the IMU executor. The node's default
callback group is assigned to the main executor. In this way, the IMU thread only
contends for shared mutexes with the main thread, not for execution time.

\subsection{Mutex Contention}
\label{sec:fr-mutex}

Even with separate executors, the IMU thread can be blocked if it attempts to
acquire a mutex held by the main thread. In the original implementation, the
\gls{lidar} callback acquires the buffer mutex before performing scan
preprocessing (point filtering, format conversion), an operation that can take
several milliseconds.

The solution is to move preprocessing outside the mutex-protected section. The
mutex is acquired only for inserting the preprocessed scan into the buffer, an
operation that takes microseconds. This reduces the mutex hold time by three orders
of magnitude.

\subsection{Path Serialization}
\label{sec:fr-path}

An additional bottleneck arises from publishing the \texttt{nav\_msgs::Path}
message, which contains the complete sequence of visited poses. If this message is
published at 200\,Hz, its size grows continuously: after 30 seconds at 200\,Hz, the
message contains 6000 poses and its serialization consumes significant CPU time on
the IMU thread.

The solution is to throttle the path message publication to 10\,Hz. The path is
visualization data, not control data, and this reduction does not affect system
functionality.

\subsection{Race Condition in Synchronization}
\label{sec:fr-sync-race}

The packet synchronization function, which pairs \gls{lidar} scans with IMU
measurements to form the \gls{ieskf} inputs, accessed the shared buffers without
mutex protection in the original implementation. With a single thread, this caused
no issues. With two threads, concurrent writing from the IMU callback and reading
from the synchronization function constitutes a data race.

The solution is to protect the synchronization function with the same mutex that
protects the buffers, ensuring mutual exclusion between write and read operations.

\section{Robustness to LiDAR Interruptions}
\label{sec:fr-stale-anchor}

\subsection{Problem: Frozen Anchor}
\label{sec:fr-frozen-anchor}

The forward propagation depends on the anchor to define the integration starting
point. If \gls{lidar} scans stop arriving (for example, due to a temporary sensor
obstruction or a packet loss in the DDS network), the anchor timestamp freezes and
$\Delta t$ grows without bound.

For large $\Delta t$, the single-step integration accumulates quadratic errors: the
accelerometer is integrated over an interval where the rotation has changed
significantly relative to $\mathbf{R}_a$, invalidating the constant-rotation
approximation of equation~\eqref{eq:fr-aworld}.

\subsection{Chained Propagation}
\label{sec:fr-chaining}

When $\Delta t$ exceeds a threshold (0.5\,s in the current configuration), the
propagation is chained: the propagated state becomes the new anchor for subsequent
IMU samples. Formally, if $\Delta t_k = t_k - t_a > 0.5$\,s:

\begin{align}
    \mathbf{p}_a &\leftarrow \mathbf{p}(t_k) \notag \\
    \mathbf{v}_a &\leftarrow \mathbf{v}(t_k)
    \label{eq:fr-chain} \\
    \mathbf{R}_a &\leftarrow \mathbf{R}(t_k) \notag \\
    t_a &\leftarrow t_k \notag
\end{align}

The biases and gravity are not updated during chaining, since their best estimate
remains that of the last \gls{ieskf}. In this way, each subsequent integration
operates on a small $\Delta t$ (the IMU period, 5\,ms), avoiding quadratic error
accumulation.

\subsection{Anchor Timestamp Guard}
\label{sec:fr-timestamp-guard}

When the processing of a \gls{lidar} scan takes longer than the inter-scan period,
the \gls{ieskf} may attempt to overwrite the anchor with a state whose timestamp is
earlier than what chained propagation has already advanced to. This would cause a
backward jump in the odometry.

To prevent this, the writing of position, velocity, rotation, and timestamp to the
anchor is conditioned on the new timestamp being later than the current one. Biases
and gravity are always updated, as they come from the \gls{ieskf} correction and are
independent of the timestamp.

\subsection{Continuity vs.\ Accuracy}
\label{sec:fr-continuity}

Chained propagation produces odometry with increasing drift during \gls{lidar}
interruptions. However, for a flight controller, drifting odometry is preferable to
no odometry at all: the controller can maintain hover with degraded feedback, but
without feedback it enters emergency mode. The drift accumulated during typical
interruptions (1--2\,s) is on the order of centimeters and is corrected when the
\gls{lidar} resumes normal operation.

\section{Chapter Summary}
\label{sec:fr-summary}

This chapter has presented the modifications required to transform a \gls{lidar}-%
inertial odometry system designed for offline mapping into one suitable for
real-time flight control. Table~\ref{tab:fr-summary} summarizes the identified
limitations and the adopted solutions.

\begin{table}[ht]
\centering
\caption{Summary of limitations and solutions for IMU-rate odometry.}
\label{tab:fr-summary}
\begin{tabular}{p{4.5cm} p{4.5cm} p{3.5cm}}
\toprule
\textbf{Limitation} & \textbf{Solution} & \textbf{Result} \\
\midrule
Odometry at 10\,Hz (LiDAR rate) &
IMU propagation from corrected anchor &
Constant 200\,Hz \\
\addlinespace
No velocities published &
$\mathbf{v}_b = \mathbf{R}^\top \mathbf{v}$,
$\boldsymbol{\omega}_b = \boldsymbol{\omega}_m - \mathbf{b}_g$ &
Full twist \\
\addlinespace
Noise in propagated velocity &
EMA filter ($\alpha = 0.05$, $f_c \approx 1.6$\,Hz) +
SLERP for orientation &
Smooth signal \\
\addlinespace
Single thread blocks IMU callback &
Separate executors + minimal mutex &
No blocking \\
\addlinespace
Odometry dies if LiDAR is interrupted &
Chained propagation with timestamp guard &
Guaranteed continuity \\
\bottomrule
\end{tabular}
\end{table}

The modifications presented do not alter the underlying estimation algorithm: the
\gls{ieskf} with point-to-plane registration and the ikd-Tree remain intact. The
changes are purely architectural and concern publishing, making them applicable to
any FAST-LIO2-based system or similar systems where the filter correction operates
at a rate lower than the IMU.
